\scanpage{154}
% archive foliation: 153
\margin{A. GROTHENDIECK.---Publications mathématiques, n°4.\quad Placard 36\quad Imp.\quad P. U. F.}
\textit{Proposition} \textbf{(6.4.12).} --- \textit{Soient $f:X\to Y$, $g:Y'\to Y$ deux morphismes; posons $X'=X\times_Y Y'$ et soit $f'=f_{(Y')}:X'\to Y'$. Soit $y'\in Y'$, $y=g(y')$; si la fibre $f^{-1}(y)$ est un schéma algébrique fini sur $\kappa(y)$, alors la fibre $f'^{-1}(y')$ est un schéma algébrique fini sur $\kappa(y')$, ayant même rang et même nombre géométrique de points que $f^{-1}(y)$.}

Compte tenu de la transitivité des fibres (3.6.5), cela résulte aussitôt de (6.4.6) et (6.4.8).

\textbf{(6.4.13)} La prop. (6.4.11) montre que les morphismes de type fini correspondent intuitivement aux « familles algébriques de variétés algébriques », les points de Y jouant le rôle de « paramètres »; ce \deleted{sont donc les seuls}\inserted{qui donne à ces} morphismes \deleted{qui aient} une signification « géométrique ». Les morphismes qui ne sont pas de type fini interviendront surtout par la suite dans les questions de « changement du préschéma de base », par exemple par localisation ou complétion.
\margin{$\times$}

\textbf{6.5. Détermination locale d'un morphisme.}

\textit{Proposition} \textbf{(6.5.1).} --- \textit{Soient X, Y deux S-préschémas, Y étant de type fini sur S; soient $x\in X$, $y\in Y$ au-dessus d'un même point $s\in S$.}

(i) \textit{Si deux S-morphismes $f=(\psi,\theta)$, $f'=(\psi',\theta')$ de X dans Y sont tels que $\psi(x)=\psi'(x)=y$, et que les $\mathcal O_s$-homomorphismes (locaux) $\theta_x^\sharp$ et $\theta_x'^\sharp$ de $\mathcal O_y$ dans $\mathcal O_x$ soient identiques, $f$ et $f'$ coïncident dans un voisinage ouvert de $x$.}

(ii) \textit{Supposons en outre S localement noethérien. Pour tout $\mathcal O_s$-homomorphisme local $\varphi:\mathcal O_y\to\mathcal O_x$, il existe un voisinage ouvert U de $x$ dans X et un S-morphisme $f=(\psi,\theta)$ de U dans Y tels que $\psi(x)=y$ et $\theta_x^\sharp=\varphi$.}

(i) La question étant locale sur S, X et Y, on peut supposer S, X, Y affines d'anneaux respectifs A, B, C, $f$ et $f'$ étant de la forme $({}^a\varphi,\widetilde\varphi)$ et $({}^a\varphi',\widetilde\varphi')$ respectivement, où $\varphi$ et $\varphi'$ sont deux A-homomorphismes de C dans B tels que $\varphi^{-1}(\mathfrak j_x)=\varphi'^{-1}(\mathfrak j_x)=\mathfrak j_y$, et les homomorphismes $\varphi_x$ et $\varphi'_x$ de $C_y$ dans $B_x$, déduits de $\varphi$ et $\varphi'$, sont identiques; on peut en outre supposer que C est une A-algèbre \textit{de type fini}. Soient $c_i$ $(1\leq i\leq n)$ des générateurs de la A-algèbre C, et posons $b_i=\varphi(c_i)$, $b'_i=\varphi'(c_i)$; par hypothèse, on a $b_i/1=b'_i/1$ dans l'anneau de fractions $B_x$ $(1\leq i\leq n)$. Cela signifie qu'il existe des éléments $s_i\in B-\mathfrak j_x$ tels que $s_i(b_i-b'_i)=0$ pour $1\leq i\leq n$, et on peut évidemment supposer tous les $s_i$ égaux à un même élément $g\in B-\mathfrak j_x$. On en conclut que l'on a $b_i/1=b'_i/1$ pour $1\leq i\leq n$ dans l'anneau de fractions $B_g$; si $i_g$ est l'homomorphisme canonique $B\to B_g$, on a par suite $i_g\circ\varphi=i_g\circ\varphi'$; donc les restrictions de $f$ et $f'$ à D $(g)$ sont identiques.

(ii) On peut se ramener à la même situation que dans (i), et supposer en outre que l'anneau A est \textit{noethérien}. Soient $c_i$ $(1\leq i\leq n)$ des générateurs de la A-algèbre C, et soit $\alpha:A[X_1,\ldots,X_n]\to C$ l'homomorphisme de l'algèbre de polynômes $A[X_1,\ldots,X_n]$ sur C transformant $X_i$ en $c_i$ pour $1\leq i\leq n$. Soit d'autre part $i_y$ l'homomorphisme canonique $C\to C_y$, et considérons l'homomorphisme composé
\diagram{\beta:\ \begin{tikzcd}
A[X_1,\ldots,X_n]\arrow[r,"\alpha"]&C\arrow[r,"i_y"]&C_y\arrow[r,"\varphi"]&B_x.
\end{tikzcd}}

